\begin{table}[htbp]
\centering
\small
\caption{V3 探索的時空間ダイナミクス（Qwen 2.5 1.5B Discovery）：直接読み出しピーク深度 $d_D$ と因果介入ピーク深度 $d_C$ の空間的解離、および下流層媒介減衰率。}
\label{tab:v3_spatiotemporal_dynamics}
\begin{tabular}{l ccccc ccc}
\toprule
 & \multicolumn{5}{c}{\textbf{Peak \& Center-of-Mass Depth}} & \multicolumn{3}{c}{\textbf{Mediated Attenuation}} \\
\cmidrule(lr){2-6} \cmidrule(lr){7-9}
\textbf{Axis} & $d_D^{\text{peak}}$ & $d_C^{\text{peak}}$ & $\Delta d^{\text{peak}}$ & $d_D^{\text{center}}$ & $d_C^{\text{center}}$ & Layer ($d$) & Atten. Ratio & 95\% CI \\
\midrule
Valence & 0.852 & 0.741 & -0.111 & 0.737 & 0.609 & L3 (0.11) & 0.87\% & [-1.77, 3.42]\% \\
Arousal & 0.852 & 0.778 & -0.074 & 0.717 & 0.756 & L3 (0.11) & 0.58\% & [0.01, 1.13]\% \\
\bottomrule
\end{tabular}
\vspace{1ex}
\begin{minipage}{\linewidth}
\footnotesize
\textbf{Note:} $\Delta d^{\text{peak}} = d_C^{\text{peak}} - d_D^{\text{peak}} < 0$ は、因果効果ピークがプローブ最適デコード深度よりも浅い（上流の）中間層に位置する「時空間解離」を示す。下流アブレーションによる減衰率はわずかであり、単一媒介経路への過度な依存がないことを示唆する。
\end{minipage}
\end{table}